% chapter3_method.tex
% 第3章 高维结构性变化检验方案的构建与验证设计
% 需要的宏包: amsmath, amssymb, amsthm, algorithm2e 或 algorithmicx

\chapter{高维结构性变化检验方案的构建与验证设计}\label{chap:method}

\section{结构性变化检验方案的构建}\label{sec:method-construction}

在处理包含多重相互依存变量的经济与金融时间序列系统时，分析框架的核心目标在于检测数据生成过程的参数是否在观测区间内保持稳定。为系统地刻画这一稳定性特征，本研究基于向量自回归（Vector Autoregression, VAR）模型，构建了一套结构性变化检验方案。该方案遵循以下推演逻辑：首先确立高维时间序列的模型设定与平稳性条件；其次，基于多元高斯似然函数推导分段似然比统计量；最后，引入基于残差重抽样的非渐近推断方法，以解决高维情形下经典渐近分布理论失效的问题。这三个模块共同构成了一个适用于不同参数结构的假设检验框架。

\subsection{检验假设的数学界定}\label{subsec:hypothesis}

\subsubsection{VAR 模型设定}

对于包含 $N$ 个观测变量的动态系统，采用滞后阶数为 $p$ 的向量自回归模型描述变量间的滞后相依关系。设 $\{Y_t\}_{t=1}^{T}$ 为 $N$ 维等间隔时间序列，其中 $Y_t \in \mathbb{R}^N$ 为第 $t$ 期的观测向量。将过去 $p$ 期的观测值按时间倒序堆叠为 $Np \times 1$ 维向量 $\mathrm{vec}(Y_{t-1}, \ldots, Y_{t-p})$，则 VAR($p$) 模型定义为：
\begin{equation}\label{eq:var-model}
    Y_t = c + \Phi \cdot \mathrm{vec}(Y_{t-1}, \ldots, Y_{t-p}) + \varepsilon_t, \quad t = p+1, \ldots, T,
\end{equation}
其中 $c \in \mathbb{R}^N$ 为截距向量；$\Phi \in \mathbb{R}^{N \times Np}$ 为系数矩阵，其第 $(i, j)$ 个元素度量了第 $j$ 个滞后变量对第 $i$ 个当前变量的边际效应；$\varepsilon_t \in \mathbb{R}^N$ 为误差向量，满足
\begin{equation}\label{eq:error-dist}
    \varepsilon_t \sim \mathcal{N}(\mathbf{0}, \Sigma), \quad \Sigma \succ 0,
\end{equation}
且不同时点的误差相互独立，即对任意 $t \neq s$，有 $\mathrm{Cov}(\varepsilon_t, \varepsilon_s) = \mathbf{0}$。

\subsubsection{平稳性条件}

为保证该系统具备良好的渐近性质与遍历性，需满足协方差平稳性条件。将 VAR($p$) 模型改写为一阶伴随形式，构造 $Np \times Np$ 维伴随矩阵：
\begin{equation}\label{eq:companion}
    \mathbf{A} =
    \begin{pmatrix}
        \Phi_1     & \Phi_2     & \cdots & \Phi_{p-1} & \Phi_p     \\
        I_N        & 0          & \cdots & 0          & 0          \\
        0          & I_N        & \cdots & 0          & 0          \\
        \vdots     & \vdots     & \ddots & \vdots     & \vdots     \\
        0          & 0          & \cdots & I_N        & 0
    \end{pmatrix},
\end{equation}
其中 $\Phi = (\Phi_1, \Phi_2, \ldots, \Phi_p)$ 为对 $\Phi$ 按 $N \times N$ 分块后的各滞后阶系数矩阵，$I_N$ 为 $N$ 阶单位矩阵。平稳性条件要求 $\mathbf{A}$ 的所有特征值的模严格小于 $1$：
\begin{equation}\label{eq:stationarity}
    \max_i |\lambda_i(\mathbf{A})| < 1.
\end{equation}
该条件保证系统对外部冲击的响应以几何速率衰减，从而使序列不会在时间维度上发散，为参数估计和统计推断提供了必要的理论基础。

\subsubsection{原假设与备择假设}

在上述模型框架下，结构性变化检验的核心任务是判断系数矩阵 $\Phi$ 是否在某一时间节点发生了改变。设 $t^*$ 为已知或预设的断点，将时间轴划分为两个不重叠的子区间 $\{p+1, \ldots, t^*\}$ 与 $\{t^*+1, \ldots, T\}$。原假设与备择假设分别定义为：
\begin{align}
    H_0 &: \Phi_1 = \Phi_2, \label{eq:h0} \\
    H_1 &: \Phi_1 \neq \Phi_2, \label{eq:h1}
\end{align}
其中 $\Phi_1$ 和 $\Phi_2$ 分别为断点前后子区间内的系数矩阵。$H_0$ 假定系统参数在整个观测期内保持不变，观测到的序列波动仅源于随机扰动 $\varepsilon_t$ 的固有方差。$H_1$ 则允许系数矩阵在断点处发生变化，只要 $\Phi_1$ 与 $\Phi_2$ 中至少存在一个元素不相等，备择假设即成立。

\subsubsection{系数矩阵的结构类型}

随着变量维度 $N$ 的增大，系数矩阵 $\Phi$ 中的自由参数数量以 $N^2 p$ 的速率增长，使得参数估计方法和检验统计量的分布特性均受到系数矩阵结构的显著影响。根据数据生成机制的差异，本文考虑以下三种典型结构：

\paragraph{低维稠密结构.}
当 $N$ 较小时，$\Phi$ 中的所有元素均为自由参数，不施加任何先验约束。所有变量之间存在密集的交叉滞后影响，可采用普通最小二乘法（OLS）直接估计。

\paragraph{中维稀疏结构.}
当 $N$ 适度增大时，变量间的传导效应往往局部化，系数矩阵中仅有少数元素为非零值。设定矩阵元素以概率 $1 - s$ 为零（其中 $s$ 为稀疏度参数），非零元素呈散点状分布。这种稀疏结构通常借助 $L_1$ 惩罚（Lasso 正则化）进行估计，通过引入绝对值惩罚项实现变量选择与降维：
\begin{equation}\label{eq:lasso}
    \hat{\Phi}_{\text{Lasso}} = \arg\min_{\Phi} \left\{ \frac{1}{2T_{\text{eff}}} \sum_{t=p+1}^{T} \|Y_t - c - \Phi \cdot \mathrm{vec}(Y_{t-1}, \ldots, Y_{t-p})\|_2^2 + \lambda \|\mathrm{vec}(\Phi)\|_1 \right\},
\end{equation}
其中 $\lambda > 0$ 为正则化参数。实际实现中采用逐方程 Lasso 回归，对 $\Phi$ 的每一行分别求解。

\paragraph{高维低秩结构.}
在宏观经济或大型金融系统中，众多变量的协同波动可能由少数潜在公共因子驱动，使得系数矩阵具有低秩性质。此时 $\Phi$ 可分解为
\begin{equation}\label{eq:lowrank}
    \Phi = U V^\top, \quad U \in \mathbb{R}^{N \times r},\; V \in \mathbb{R}^{Np \times r},
\end{equation}
其中 $r \ll \min(N, Np)$ 为矩阵的真实秩。估计方法采用降秩回归（Reduced-Rank Regression, RRR）：首先对含截距的设计矩阵 $X_{\text{full}}$ 执行 OLS 估计得到 $\hat{B}_{\text{OLS}}$，随后对拟合值矩阵 $\hat{Y} = X_{\text{full}} \hat{B}_{\text{OLS}}$ 进行奇异值分解，取前 $r$ 个右奇异向量 $V_r$，通过投影
\begin{equation}\label{eq:rrr-proj}
    \hat{B}_{\text{RRR}} = \hat{B}_{\text{OLS}} \cdot V_r V_r^\top
\end{equation}
获得最优的秩 $r$ 逼近。RRR 最小化的是预测误差 $\|Y - X\Phi\|_F^2$，是低秩 VAR 估计的经典方法。

上述三种结构虽未改变假设检验关于 $\Phi_1 = \Phi_2$ 的本质表述，但直接影响了参数估计的求解路径和检验统计量的有限样本分布特性，这要求检验方案在统计量构造和推断方法上均需适应不同的参数结构。


\subsection{分段似然比检验统计量的构造}\label{subsec:lr-statistic}

在明确了模型的假设空间与参数矩阵结构后，需将样本中蕴含的结构变化信息压缩为标量检验统计量。基于分段最大似然估计原理的似然比（Likelihood Ratio, LR）统计量是衡量断点前后参数差异的基本工具。其核心逻辑是比较数据在原假设约束空间与备择假设非约束空间下的最大对数似然值，通过二者之差的大小判断参数同质性约束所导致的拟合损失程度。

\subsubsection{集中对数似然函数}

基于误差项的多元高斯分布假定，有效观测样本\footnote{由于 VAR 方程的递推需要前 $p$ 期作为初始值，有效样本量为 $T_{\text{eff}} = T - p$。} 的对数似然函数为
\begin{equation}\label{eq:loglik-full}
    \ell(\Phi, c, \Sigma) = -\frac{T_{\text{eff}}}{2} \left[ N \ln(2\pi) + \ln|\Sigma| \right] - \frac{1}{2} \sum_{t=p+1}^{T} (Y_t - \mu_t)^\top \Sigma^{-1} (Y_t - \mu_t),
\end{equation}
其中 $\mu_t = c + \Phi \cdot \mathrm{vec}(Y_{t-1}, \ldots, Y_{t-p})$ 为条件均值，$T_{\text{eff}} = T - p$ 为有效样本量。

对 $\Sigma$ 求一阶导数并令其为零，可得残差协方差矩阵的最大似然估计量：
\begin{equation}\label{eq:sigma-mle}
    \hat{\Sigma} = \frac{1}{T_{\text{eff}}} \sum_{t=p+1}^{T} \hat{\varepsilon}_t \hat{\varepsilon}_t^\top,
\end{equation}
其中 $\hat{\varepsilon}_t = Y_t - \hat{c} - \hat{\Phi} \cdot \mathrm{vec}(Y_{t-1}, \ldots, Y_{t-p})$ 为拟合残差。将式~\eqref{eq:sigma-mle} 代入式~\eqref{eq:loglik-full}，二次型部分经矩阵迹运算化简为常数 $N$，得到集中对数似然函数：
\begin{equation}\label{eq:concentrated-loglik}
    \ell(\hat{\Phi}, \hat{c}) = -\frac{T_{\text{eff}}}{2} \left[ N \ln(2\pi) + \ln|\hat{\Sigma}| + N \right].
\end{equation}
由于在似然比的差分运算中，与 $\Phi$ 无关的常数项 $N\ln(2\pi) + N$ 将被抵消，因此检验统计量的构造仅依赖于
\begin{equation}\label{eq:loglik-core}
    \ell \propto -\frac{T_{\text{eff}}}{2} \ln|\hat{\Sigma}|.
\end{equation}

\subsubsection{约束与非约束模型下的似然计算}

\paragraph{原假设 $H_0$（约束模型）.}
在 $H_0$ 下，整个有效样本被视为由同一组参数驱动的同质系统。利用全部 $T_{\text{eff}}$ 个有效观测值，根据所选的估计方法（OLS、Lasso 或 RRR）执行全局参数拟合，获得全局系数估计量 $\hat{\Phi}^{(0)}$、截距估计量 $\hat{c}^{(0)}$ 以及残差协方差估计矩阵 $\hat{\Sigma}^{(0)}$。将 $\hat{\Sigma}^{(0)}$ 代入式~\eqref{eq:concentrated-loglik}，即得约束模型下的最大对数似然值 $\ell_0$。

\paragraph{备择假设 $H_1$（非约束模型）.}
在 $H_1$ 下，参数允许在断点 $t^*$ 处发生变化。将全样本以 $t^*$ 为界拆分为两个子样本，分别独立拟合 VAR 模型。具体地，子样本 1 包含观测值 $Y_1, \ldots, Y_{t^*}$，有效样本量为 $T_1 = t^* - p$；子样本 2 包含观测值 $Y_{t^*-p+1}, \ldots, Y_T$，有效样本量为 $T_2 = T - t^*$\footnote{子样本 2 的起始点回溯 $p$ 期以提供 VAR 方程所需的初始滞后值，从而保证两个子样本的有效观测总量 $T_1 + T_2 = T - p = T_{\text{eff}}$ 与全样本一致。}。对两个子样本分别估计得到 $(\hat{\Phi}_1^{(1)}, \hat{\Sigma}_1^{(1)})$ 和 $(\hat{\Phi}_2^{(1)}, \hat{\Sigma}_2^{(1)})$，非约束模型下的总对数似然值为两个子样本对数似然值之和：
\begin{equation}\label{eq:loglik-h1}
    \ell_1 = \ell(\hat{\Phi}_1^{(1)}) + \ell(\hat{\Phi}_2^{(1)}) = -\frac{T_1}{2}\left[N\ln(2\pi) + \ln|\hat{\Sigma}_1^{(1)}| + N\right] - \frac{T_2}{2}\left[N\ln(2\pi) + \ln|\hat{\Sigma}_2^{(1)}| + N\right].
\end{equation}

\subsubsection{似然比统计量}

分段似然比检验统计量定义为非约束与约束模型最大对数似然值之差的两倍：
\begin{equation}\label{eq:lr-statistic}
    \mathrm{LR}(t^*) = 2\left[\ell_1 - \ell_0\right] = T_{\text{eff}} \ln|\hat{\Sigma}^{(0)}| - T_1 \ln|\hat{\Sigma}_1^{(1)}| - T_2 \ln|\hat{\Sigma}_2^{(1)}|.
\end{equation}
该统计量具有直观的统计含义：当数据生成机制不存在结构性变化时，全样本拟合与分段拟合的残差协方差矩阵应接近一致，$\mathrm{LR}(t^*)$ 的值较小；反之，若系统在断点处发生参数变化，全样本约束拟合将产生模型误设，残差中将积累本应由系数变化解释的动态信息，导致 $\hat{\Sigma}^{(0)}$ 的行列式相对于 $\hat{\Sigma}_1^{(1)}$ 和 $\hat{\Sigma}_2^{(1)}$ 显著增大，从而使 $\mathrm{LR}(t^*)$ 取较大值。


\subsection{基于重抽样的 \texorpdfstring{$p$}{p} 值计算与统计决策准则}\label{subsec:bootstrap}

\subsubsection{渐近分布失效与 Bootstrap 方法的引入}

在低维情形下，似然比统计量在 $H_0$ 为真时渐近服从自由度为 $N^2 p$ 的 $\chi^2$ 分布。然而，在高维 VAR 模型中，待估参数数量以 $O(N^2 p)$ 的速率增长，当 $N$ 较大时，参数与样本量之比不再接近零，$\chi^2$ 渐近近似失效。具体原因包括：(1) 为克服参数爆炸而引入的 Lasso 惩罚或降秩约束改变了参数空间的几何性质，使得估计量不再具备渐近正态性；(2) 正则化参数的数据驱动选择引入了额外的随机性，进一步扭曲统计量的抽样分布。这些因素使得似然比统计量的真实有限样本分布偏离 $\chi^2$ 分布。若仍依赖渐近临界值进行推断，将产生严重的水平扭曲，即检验以远高于名义水平的频率错误拒绝 $H_0$。

为克服上述问题，本方案采用基于残差的 Bootstrap 重抽样方法。Bootstrap 通过对样本内部信息的随机重组，以经验分布逼近检验统计量在有限样本下的真实分布，无需依赖渐近理论的正则条件。对于 VAR 等具有时间依赖结构的动态系统，不能直接对观测向量进行简单的随机抽样，否则会破坏序列的自相关结构。残差 Bootstrap 的策略是保留 $H_0$ 下拟合模型所确定的系统动态方程，仅对模型残差进行随机重抽样，由此生成在 $H_0$ 下的伪时间序列。

\subsubsection{残差 Bootstrap 重抽样过程}

基于残差的 Bootstrap 推断包含以下步骤：

\paragraph{步骤 1：原始数据拟合与基准统计量计算.}
在 $H_0$ 下对全样本拟合 VAR($p$) 模型，提取全局系数估计 $\hat{\Phi}^{(0)}$、截距估计 $\hat{c}^{(0)}$ 及拟合残差 $\hat{\varepsilon}_t$（$t = p+1, \ldots, T$）。同时按式~\eqref{eq:lr-statistic} 计算原始数据的似然比统计量 $\mathrm{LR}_{\text{obs}}$。

\paragraph{步骤 2：残差中心化.}
为消除有限样本下残差均值的微小偏移，计算中心化残差：
\begin{equation}\label{eq:center-resid}
    \tilde{\varepsilon}_t = \hat{\varepsilon}_t - \bar{\varepsilon}, \quad \bar{\varepsilon} = \frac{1}{T_{\text{eff}}} \sum_{t=p+1}^{T} \hat{\varepsilon}_t.
\end{equation}

\paragraph{步骤 3：Bootstrap 循环.}
设重抽样次数为 $B$。对 $b = 1, \ldots, B$，执行以下操作：

\begin{enumerate}
    \item \textbf{残差重抽样.} 从中心化残差集合 $\{\tilde{\varepsilon}_{p+1}, \ldots, \tilde{\varepsilon}_T\}$ 中以均匀概率有放回地抽取 $T_{\text{eff}}$ 个向量，构成伪残差序列 $\{\tilde{\varepsilon}_t^{*b}\}_{t=p+1}^{T}$。

    \item \textbf{伪序列生成.} 保留原始数据的前 $p$ 个观测值作为初始值 $Y_{1:p}^{*b} = Y_{1:p}$，从 $t = p+1$ 起按下式逐期递推：
    \begin{equation}\label{eq:bootstrap-recursion}
        Y_t^{*b} = \hat{c}^{(0)} + \hat{\Phi}^{(0)} \cdot \mathrm{vec}(Y_{t-1}^{*b}, \ldots, Y_{t-p}^{*b}) + \tilde{\varepsilon}_t^{*b}.
    \end{equation}
    该递推保证每条伪序列在数据生成机制上服从 $H_0$ 下的全局恒定参数结构，同时保留原始数据的波动特征。

    \item \textbf{伪统计量计算.} 对伪序列 $Y^{*b}$ 执行与原始数据完全一致的检验流程——按所选估计方法分别进行全样本约束拟合和断点分段非约束拟合，计算伪似然比统计量 $\mathrm{LR}^{*b}$。在拟合过程中，Lasso 的逐方程惩罚回归或 RRR 降秩投影均需一致地重新执行。
\end{enumerate}

\paragraph{步骤 4：经验 $p$ 值计算与统计决策.}
经验 $p$ 值定义为 Bootstrap 伪统计量中不小于观测统计量的比例：
\begin{equation}\label{eq:p-value}
    \hat{p} = \frac{1}{B} \sum_{b=1}^{B} \mathbf{1}\{\mathrm{LR}^{*b} \geq \mathrm{LR}_{\text{obs}}\}.
\end{equation}
该值度量了在 $H_0$ 为真的条件下，仅由随机波动产生的统计量超过观测值的频率。给定名义显著性水平 $\alpha$，决策准则为：
\begin{equation}\label{eq:decision}
    \text{若 } \hat{p} \leq \alpha, \text{ 则拒绝 } H_0; \quad \text{若 } \hat{p} > \alpha, \text{ 则不拒绝 } H_0.
\end{equation}

\subsubsection{检验程序伪代码}

以上流程的完整执行逻辑总结为算法~\ref{alg:bootstrap-test}。

\begin{algorithm}[htbp]
\caption{基于残差 Bootstrap 的高维 VAR 结构性变化检验}\label{alg:bootstrap-test}
\KwIn{时间序列 $Y \in \mathbb{R}^{T \times N}$，断点 $t^*$，滞后阶数 $p$，重抽样次数 $B$，显著性水平 $\alpha$}
\KwOut{检验决策（拒绝或不拒绝 $H_0$）}

\tcp{步骤 1：原始数据拟合与基准统计量计算}
对全样本拟合约束 VAR($p$) 模型，提取 $\hat{\Phi}^{(0)}$、$\hat{c}^{(0)}$ 与残差 $\{\hat{\varepsilon}_t\}$\;
以 $t^*$ 为断点分段拟合非约束 VAR($p$) 模型\;
计算观测似然比统计量 $\mathrm{LR}_{\text{obs}}$\;

\tcp{步骤 2：残差中心化}
$\bar{\varepsilon} \leftarrow \frac{1}{T_{\text{eff}}} \sum_{t=p+1}^{T} \hat{\varepsilon}_t$\;
$\tilde{\varepsilon}_t \leftarrow \hat{\varepsilon}_t - \bar{\varepsilon}$，$t = p+1, \ldots, T$\;

\tcp{步骤 3：Bootstrap 循环}
$\mathrm{count} \leftarrow 0$\;
\For{$b = 1$ \KwTo $B$}{
    从 $\{\tilde{\varepsilon}_{p+1}, \ldots, \tilde{\varepsilon}_T\}$ 中有放回抽样 $T_{\text{eff}}$ 个向量 $\rightarrow$ $\{\tilde{\varepsilon}_t^{*b}\}$\;
    $Y_{1:p}^{*b} \leftarrow Y_{1:p}$\;
    \For{$t = p+1$ \KwTo $T$}{
        $Y_t^{*b} \leftarrow \hat{c}^{(0)} + \hat{\Phi}^{(0)} \cdot \mathrm{vec}(Y_{t-1}^{*b}, \ldots, Y_{t-p}^{*b}) + \tilde{\varepsilon}_t^{*b}$\;
    }
    对 $Y^{*b}$ 执行全样本约束拟合与断点 $t^*$ 分段非约束拟合\;
    \tcp{拟合方法需与原始数据一致（Lasso 惩罚或 RRR 降秩投影）}
    计算伪统计量 $\mathrm{LR}^{*b}$\;
    \If{$\mathrm{LR}^{*b} \geq \mathrm{LR}_{\emph{obs}}$}{
        $\mathrm{count} \leftarrow \mathrm{count} + 1$\;
    }
}

\tcp{步骤 4：经验 $p$ 值与决策}
$\hat{p} \leftarrow \mathrm{count} / B$\;
\eIf{$\hat{p} \leq \alpha$}{
    \Return 拒绝 $H_0$：存在结构性变化\;
}{
    \Return 不拒绝 $H_0$：参数保持稳定\;
}
\end{algorithm}

该检验方案在理论上避免了对高维渐近分布的依赖，在计算框架上对稠密、稀疏和低秩等不同参数结构均具有适应性，为高维 VAR 模型下的结构性变化推断提供了可行的计算基础。
